\chapter{Trabalhos Relacionados}
\label{sec:relacionados}


Este capítulo apresenta os principais trabalhos sobre o uso de Modelos de Linguagem de Grande Porte (LLMs) em contextos de avaliação, certificação e apoio ao aprendizado, com o objetivo de posicionar o presente estudo em relação ao estado da arte.


\section{LLMs em Avaliações Educacionais}
\label{subsec:llm_educacao}

\textbf{Ideias a abordar:}
\begin{itemize}
    \item Estudos que avaliam LLMs em exames acadêmicos:
    \begin{itemize}
        \item ENEM, GMAT, exames universitários
        \item Provas de áreas como medicina, computação e negócios
    \end{itemize}
    \item Principais resultados:
    \begin{itemize}
        \item Desempenho comparável ou superior a humanos em questões textuais
        \item Melhor desempenho em questões objetivas e factuais
    \end{itemize}
    \item Limitações identificadas:
    \begin{itemize}
        \item Dificuldade em questões que exigem raciocínio profundo
        \item Problemas com ambiguidade e interpretação
        \item Sensibilidade ao idioma e formulação da questão
    \end{itemize}
    \item Lacuna:
    \begin{itemize}
        \item Pouca exploração de exames profissionais complexos
    \end{itemize}
\end{itemize}



\section{LLMs em Certificações Profissionais}
\label{subsec:llm_certificacao}

\textbf{Ideias a abordar:}
\begin{itemize}
    \item Trabalhos que avaliam LLMs em certificações:
    \begin{itemize}
        \item Scrum (PSM, PSPO)
        \item Certificações técnicas em TI
    \end{itemize}
    \item Características desses estudos:
    \begin{itemize}
        \item Uso de datasets estruturados de questões
        \item Avaliação de acurácia e consistência
    \end{itemize}
    \item Resultados típicos:
    \begin{itemize}
        \item Alto desempenho em conteúdo normativo
        \item Dificuldades em questões interpretativas ou ambíguas
    \end{itemize}
    \item Limitações:
    \begin{itemize}
        \item Domínios mais restritos e bem definidos (ex.: Scrum Guide)
        \item Pouca variabilidade contextual
    \end{itemize}
    \item Lacuna:
    \begin{itemize}
        \item Ausência de estudos no contexto PMP
        \item Falta de análise em cenários híbridos e situacionais
    \end{itemize}
\end{itemize}




\section{Influência de Técnicas de Prompting}
\label{subsec:prompt_related}

\textbf{Ideias a abordar:}
\begin{itemize}
    \item Estudos sobre \textit{prompt engineering}:
    \begin{itemize}
        \item Zero-shot
        \item Few-shot
        \item Chain-of-Thought (CoT)
        \item Prompt com grounding (uso de fontes)
    \end{itemize}
    \item Evidências empíricas:
    \begin{itemize}
        \item Prompts estruturados tendem a melhorar a acurácia
        \item CoT ajuda em tarefas de raciocínio
        \item Grounding reduz alucinações
    \end{itemize}
    \item Limitações:
    \begin{itemize}
        \item Ganhos dependem da tarefa
        \item Nem sempre melhora em cenários complexos
    \end{itemize}
    \item Relevância para este trabalho:
    \begin{itemize}
        \item PMP envolve interpretação e decisão → forte dependência de prompting
    \end{itemize}
\end{itemize}


\section{Síntese dos Trabalhos e Lacuna de Pesquisa}
\label{subsec:sintese_gap}

\textbf{Ideias a abordar:}
\begin{itemize}
    \item Síntese dos principais achados:
    \begin{itemize}
        \item LLMs performam bem em tarefas estruturadas e factuais
        \item desempenho varia com o tipo de questão e prompt
    \end{itemize}
    \item Limitações consolidadas:
    \begin{itemize}
        \item dificuldade em tarefas situacionais
        \item pouca robustez em julgamento contextual
    \end{itemize}
    \item Gap explícito:
    \begin{itemize}
        \item ausência de estudos no contexto do PMP atual
        \item falta de análise integrada de:
        \begin{itemize}
            \item domínios do exame
            \item tipos de questão
            \item estratégias de prompting
        \end{itemize}
    \end{itemize}
    \item Posicionamento do trabalho:
    \begin{itemize}
        \item este estudo busca preencher essa lacuna
    \end{itemize}
    \item Transição para o Capítulo de Metodologia:
    \begin{itemize}
        \item introduzir que o próximo capítulo descreve o desenho experimental adotado
    \end{itemize}
\end{itemize}


\section{Considerações Finais do Capítulo}
\label{subsec:sintese_gap}